% =====================================================================
%  CARNET D'ENTRAINEMENT — FICHE 11 — THEME 3
%  Tuto : Équilibrer une équation d'oxydo-réduction
% =====================================================================
\documentclass[11pt,a4paper]{article}
\usepackage[utf8]{inputenc}
\usepackage[T1]{fontenc}
\usepackage{lmodern}
\usepackage[french]{babel}
\usepackage{amsmath,amssymb}
\usepackage[version=4]{mhchem}
\usepackage{siunitx}
\sisetup{output-decimal-marker={,},per-mode=symbol}
\usepackage{xcolor}
\usepackage{array}
\usepackage{booktabs}
\usepackage{enumitem}
\usepackage[most]{tcolorbox}
\usepackage[a4paper,margin=2.0cm,headheight=26pt,headsep=10pt]{geometry}
\usepackage{fancyhdr}
\usepackage{titlesec}
\usepackage[hidelinks]{hyperref}

\definecolor{primary}{RGB}{146,95,160}
\definecolor{primarydark}{RGB}{92,52,104}
\definecolor{defcol}{RGB}{0,114,178}
\definecolor{thmcol}{RGB}{0,140,120}
\definecolor{excol}{RGB}{0,135,90}
\definecolor{attncol}{RGB}{200,40,55}
\definecolor{methcol}{RGB}{90,90,100}

\newcommand{\NO}{\ensuremath{\mathrm{N.O.}}}

\tcbset{boxbase/.style={enhanced,breakable,boxrule=0.9pt,arc=2.5pt,
   left=3mm,right=3mm,top=2.3mm,bottom=2.3mm,fonttitle=\bfseries,coltitle=white,
   toptitle=0.6mm,bottomtitle=0.6mm}}
\newtcolorbox{defi}[1][]{boxbase,colback=defcol!5,colframe=defcol,title={\textsc{Définition}\ifx\relax#1\relax\relax\else\textmd{ --- \itshape #1}\fi}}
\newtcolorbox{regle}[1][]{boxbase,colback=thmcol!6,colframe=thmcol,title={\textsc{Règle}\ifx\relax#1\relax\relax\else\textmd{ --- \itshape #1}\fi}}
\newtcolorbox{exemple}[1][]{boxbase,colback=excol!5,colframe=excol,title={\textsc{Exemple}\ifx\relax#1\relax\relax\else\textmd{ : \itshape #1}\fi}}
\newtcolorbox{attention}[1][]{boxbase,colback=attncol!6,colframe=attncol,title={\textsc{Attention}\ifx\relax#1\relax\relax\else\textmd{ : \itshape #1}\fi}}
\newtcolorbox{methode}[1][]{boxbase,colback=methcol!7,colframe=methcol,sharp corners,title={\textsc{Méthode}\ifx\relax#1\relax\relax\else\textmd{ : \itshape #1}\fi}}

\pagestyle{fancy}\fancyhf{}
\fancyhead[L]{\small\textcolor{primarydark}{\textbf{Fiche 11} — Tuto Thème 3}}
\fancyhead[R]{\small\textcolor{primarydark}{\textsc{Équilibrer une équation redox}}}
\fancyfoot[C]{\small\thepage}
\renewcommand{\headrulewidth}{0.8pt}
\renewcommand{\headrule}{\hbox to\headwidth{\color{primary}\leaders\hrule height \headrulewidth\hfill}}
\titleformat{\section}{\large\bfseries\color{primarydark}}{\thesection}{0.6em}{}

\begin{document}
\begin{center}
{\LARGE\bfseries\color{primarydark}Thème 3 — Équilibrer une équation d'oxydo-réduction}\\[3pt]
{\color{primary}\rule{9.5cm}{1.2pt}}\\[4pt]
{\small\itshape À lire avant les exercices — la méthode des demi-équations, en milieu acide et basique.}
\end{center}

\vspace{-2pt}
\section*{1. La méthode en 6 étapes}
\begin{methode}[équilibrer une équation redox]
\begin{enumerate}[label=\textbf{\arabic*.},leftmargin=2.2em,topsep=2pt,itemsep=1.5pt]
  \item \textbf{Repérer les deux couples} : les deux atomes dont le \NO\ varie $\to$ deux demi-équations.
  \item \textbf{Compter les électrons} : $n=\NO(\text{ox})-\NO(\text{red})$, placés \textbf{du côté de l'oxydant}.
  \item \textbf{Conserver l'élément central} : s'il y a 2 atomes d'un côté et 1 de l'autre, multiplier
        \emph{aussi} les électrons.
  \item \textbf{Ajuster le milieu} (voir §2).
  \item \textbf{Contrôler} : mêmes atomes \emph{et} mêmes charges de chaque côté.
  \item \textbf{Équation globale} : multiplier pour égaliser les électrons, additionner, simplifier.
\end{enumerate}
\end{methode}

\section*{2. Ajuster le milieu}
\begin{regle}[milieu acide]
Ordre : \textbf{(1)} équilibrer l'\textbf{oxygène} avec des \ce{H2O} du côté qui en manque ;
\textbf{(2)} équilibrer l'\textbf{hydrogène} avec des \ce{H+} de l'autre côté ;
\textbf{(3)} ajouter les \ce{e^-} pour équilibrer la charge.
\end{regle}

\begin{regle}[milieu basique]
On équilibre \textbf{d'abord comme en milieu acide}, puis on ajoute autant de \ce{OH^-} des
\textbf{deux côtés} qu'il y a de \ce{H+}. Chaque $\ce{H+ + OH^-}$ devient une \ce{H2O}, et l'on
simplifie les \ce{H2O} en excès. \emph{(Variante : certains énoncés utilisent \ce{H3O+} au lieu de
\ce{H+} — un \ce{H+} « habillé » d'une \ce{H2O}.)}
\end{regle}

\section*{3. Exemple complet — milieu acide}
\begin{exemple}[\ce{NO3^- + Fe^{2+} -> NO + Fe^{3+}} en milieu acide]
\textbf{Couples \& électrons :} \ce{NO3^-}/\ce{NO} (N : $+\mathrm{V}\to+\mathrm{II}$, $n=3$) et
\ce{Fe^{3+}}/\ce{Fe^{2+}} (Fe : $+\mathrm{II}\to+\mathrm{III}$, $n=1$).\\[2pt]
\textbf{Demi-équations équilibrées (acide) :}
\[
\ce{NO3^- + 4 H+ + 3 e^- -> NO + 2 H2O} \qquad ; \qquad \ce{Fe^{2+} -> Fe^{3+} + e^-}.
\]
\textbf{Global :} on multiplie la seconde par 3 (pour 3 e$^-$), puis on additionne :
\[
\boxed{\ce{NO3^- + 3 Fe^{2+} + 4 H+ -> NO + 3 Fe^{3+} + 2 H2O}}
\]
\textbf{Contrôle :} charges $-1+6+4=+9$ à gauche $=+9$ à droite ($3\times(+3)$) ; atomes conservés. \checkmark
\end{exemple}

\section*{4. Exemple — milieu basique (demi-équation)}
\begin{exemple}[\ce{MnO4^- -> MnO2} en milieu basique]
Mn : $+\mathrm{VII}\to+\mathrm{IV}$, $n=3$. En \textbf{acide} : $\ce{MnO4^- + 4 H+ + 3 e^- -> MnO2 + 2 H2O}$.
On ajoute 4 \ce{OH^-} des deux côtés, les $4(\ce{H+}+\ce{OH^-})$ font 4 \ce{H2O}, on simplifie 2 \ce{H2O} :
\[
\ce{MnO4^- + 2 H2O + 3 e^- -> MnO2 + 4 OH^-}.
\]
Contrôle des charges : $-1-3=-4$ à gauche $=-4$ à droite. \checkmark
\end{exemple}

\begin{attention}[l'erreur qui coûte des points]
Quand l'oxydant contient \textbf{2} atomes de l'élément central (ex. \ce{Cr2O7^{2-} -> 2 Cr^{3+}}),
il faut multiplier le produit \emph{et} les électrons : $\Delta\NO=3$ par Cr, mais $2\times3=6$
électrons au total. Ne jamais oublier ce facteur !
\end{attention}
\end{document}
